\documentclass[../report_main.tex]{subfiles}
\begin{document}

The rapid evolution of digital education has shifted the pedagogical focus from traditional classroom management toward integrated platforms that facilitate remote learning and automated assessments. Within this paradigm, the underlying database serves as the critical foundation for ensuring data consistency, operational security, and system performance. An effective \textit{Online Examination and Question Bank Management Platform} demands a sophisticated relational architecture capable of handling diverse data structures, complex hierarchical dependencies, and real-time result aggregation. This project, conducted for the Database Lab course, investigates the structural design and practical implementation of such an ecosystem utilizing PostgreSQL.

To build such an ecosystem, one must address several fundamental problems. Architecting a database for dynamic online examinations introduces unique technical challenges that conventional management systems frequently fail to address. Primarily, the question repository must be sufficiently versatile to support heterogeneous formats-including \gls{MC}, \gls{TF}, \gls{SA}, and \gls{OE} items-while simultaneously accommodating multimedia integration and hierarchical parent-child relationships for passage-based inquiries (\gls{ST}). Furthermore, the system must enforce strict referential integrity to ensure that student submissions are accurately mapped to specific temporal states of an examination. Finally, authorization hierarchies must be structurally embedded to prevent unauthorized access to academic records, ensuring that educators can exclusively retrieve data pertaining to their designated classes. Existing legacy solutions often exhibit data redundancy and schema inflexibility when attempting to satisfy these concurrent requirements \cite{db_project_guidelines}.

Driven by these challenges, the primary objective of this research is to architect and deploy a normalized relational database schema that comprehensively supports the lifecycle of the examination process. This objective is subdivided into four fundamental goals. First, the project proposes a robust \gls{IAM} model to explicitly delineate administrative, educator, and student roles. Second, it establishes an organized academic schema where interactions between students and teachers are mediated through defined classroom entities. Third, the study develops a polymorphic Question Bank engineered to support complex metadata, difficulty stratification, and randomized components. Lastly, the system incorporates a flexible Assessment Engine that leverages semi-structured data types (JSONB) to manage dynamic scoring configurations, thereby facilitating automated grading and analytical reporting.

To realize these objectives, the developmental methodology relies on robust relational technologies. The proposed architecture is deployed on the PostgreSQL \cite{postgresql2023} relational database management system, comprising interconnected tables designed to minimize data anomaly and maximize query efficiency. The structural integrity of the database is maintained through a rigorous combination of primary, foreign, and unique key constraints based on principles established in \gls{HUST}'s internal curricula \cite{hust_db_slides}. To address modern application requirements, the design integrates Universally Unique Identifiers (UUIDs) for external exposure and precise TIMESTAMP logging for auditability. The developmental methodology strictly adheres to database normalization principles, ensuring theoretical soundness and operational scalability \cite{uet_db_slides}.

To guide the reader through the subsequent chapters, the remainder of this document is structured to provide a systematic overview of the database engineering process. Chapter 2 details the Entity-Relationship mapping and provides a comprehensive specification of the physical schema and indexing strategies. Chapter 3 examines the implementation of complex business logic through SQL, presenting advanced queries, procedural functions, and triggers. Chapter 4 illustrates the integration of the database with a prototype application interface. Finally, Chapter 5 synthesizes the project outcomes, delineates implementation constraints, and proposes avenues for future architectural enhancements.

\vspace{0.5cm}
\begin{tcolorbox}[colback=white, colframe=black!70, boxrule=0.5pt, arc=2pt, fonttitle=\bfseries, title=Project Resources]
The complete application source code and database backup have been made publicly available for the evaluation committee's reference:

\begin{itemize}
    \item \textbf{Application Source Code (GitHub):} \url{https://github.com/PhungMinhPhuc/question-dataset}
    \item \textbf{Database Backup File (Google Drive):} \url{https://drive.google.com/drive/folders/1SPMrp4qBzDp1KWkO36wYqUqM9LMyFtr_}
\end{itemize}
\end{tcolorbox}

\end{document}